\label{m1:sec:tools}\label{m2:sec:tools}
This section collects the distributional tools used in both
multiplicity-one proofs. We recall the Gelfand--Kazhdan criterion
and generalized Harish--Chandra descent, together with the
localization and homogeneity arguments needed for the normal spaces.

Throughout the paper $F$ is a non-Archimedean local field of characteristic
zero, and $\Fb$ is an algebraic closure of $F$. An $l$-space is a locally
compact totally disconnected Hausdorff space; an $l$-group is a topological
group with this property. Write $\mathcal S(X)=C_c^\infty(X,\mathbf C)$ and
$\mathcal D(X)=\mathcal S(X)^*$ for its Schwartz functions and distributions. For an
$l$-group $L$ acting on $X$, our conventions are
\begin{equation}\label{m1:eq:distribution-conventions}
 (l\cdot f)(x)=f(l^{-1}x),\qquad (lD)(f)=D(f\circ l),\qquad
 \mathcal D(X)^{L,\chi}=\{D:lD=\chi(l)D\}.
\end{equation}
In particular, an invertible linear map $l$ acts on additive Haar measure
by $|\det l|^{-1}$. We omit $\chi$ when it is trivial.

\begin{proposition}[Open and closed support]\label{m2:tools:support}
For a closed subset $Z\subset X$ and its open complement $U$, there
is an exact sequence
\[
 0\longrightarrow\mathcal D(Z)\longrightarrow\mathcal D(X)
   \longrightarrow\mathcal D(U)\longrightarrow0.
\]
The first arrow identifies $\mathcal D(Z)$ with distributions supported
in $Z$, and the second arrow is restriction to the open set $U$.
If a group preserves $Z$, these maps are equivariant. In particular,
vanishing of a fixed equivariance space on both $U$ and $Z$ implies
its vanishing on $X$.
\end{proposition}

\begin{proof}
Compact-open extension of locally constant functions gives the exact
sequence
\[
 0\longrightarrow\mathcal S(U)\longrightarrow\mathcal S(X)
  \longrightarrow\mathcal S(Z)\longrightarrow0.
\]
Dualizing proves the assertion; taking equivariants preserves its
injection and kernel. In particular, a distribution supported at one
point is a multiple of its Dirac distribution.
\end{proof}

\subsection{The Gelfand--Kazhdan criterion and descent}
\label{m1:subsec:GK}
Let $G$ be reductive and $H\subset G$ a closed reductive subgroup. Suppose
that $\tau$ is an involutive anti-automorphism with $\tau(H)=H$. The action
of $H\times H$ is $(h_1,h_2)g=h_1gh_2^{-1}$.

\begin{proposition}[Gelfand--Kazhdan criterion]\label{m1:prop:GK}\label{m2:tools:gk}
If every $H\times H$-invariant distribution on $G$ is $\tau$-invariant,
then every irreducible admissible smooth complex representation $\pi$ of
$G$ satisfies
\[
 \dim\Hom_H(\pi,\Co)\,\dim\Hom_H(\pi^\vee,\Co)\leq1.
\]
\end{proposition}
This is \cite[Lemma~4.2]{Prasad90}; see also
\cite[Remark~2.4.6]{AGSGelfand}. The invariant measure on $G/H$ required in
Prasad's formulation exists because $G$ and $H$ are unimodular.
Contragredient duality will identify the two multiplicities for each block pair.

Whenever an involution $\tau$ normalizes an action of $L$, write
$\widetilde L=L\rtimes\langle\tau\rangle$ and let $\varepsilon$ be its
sign character: $\varepsilon|_L=1$ and $\varepsilon(\tau)=-1$.
Write $\mathcal D(X)^{L,\tau}$ for distributions fixed by both $L$ and $\tau$.
Transpose invariance of $\mathcal D(X)^L$ is equivalent to
$\mathcal D(X)^{\widetilde L,\varepsilon}=0$, by taking the odd part of a
distribution. We retain this convention for all extended groups.

For an algebraic action of a reductive group $L$ on a smooth affine
variety $X$, write $L_x$ for the stabilizer of $x$ and put
$N_x=T_xX/T_x(Lx)$, the normal space to the orbit in the ambient variety $X$.
We write $\mathcal D(Y)$ for $\mathcal D(Y(F))$ when $Y$ is an $F$-variety.
A point $x\in X(F)$ is semisimple if its $L(F)$-orbit
is analytically closed, equivalently its geometric orbit is Zariski
closed \cite[Theorem~2.3.8]{AGDescent}.

\begin{theorem}[Generalized Harish--Chandra descent]
\label{m1:thm:HC-descent}\label{m2:tools:descent}\label{m1:subsec:HC-descent}
Let $L$ be reductive, possibly disconnected, acting on a smooth affine
$F$-variety $X$, and let $\chi$ be a character of $L(F)$. If
$\mathcal D(N_x(F))^{L_x(F),\chi}=0$ for every semisimple $x\in X(F)$,
then $\mathcal D(X(F))^{L(F),\chi}=0$.
\end{theorem}
We use \cite[Theorem~3.1.1]{AGDescent}, with the convention on disconnected
reductive groups of that paper.

Suppose now that $\tau$ and the automorphism $\alpha$ of $L$ that it
induces are algebraic over $F$, and that every closed rational orbit is
$\tau$-stable. Choose $g_x\in L(F)$ with $g_x\tau(x)=x$ and set
$\tau_x=g_x\tau$. Then
\begin{equation}\label{m1:eq:corrected-extension}
 \widetilde L_x=L_x(F)\sqcup L_x(F)\tau_x,
 \qquad \tau_x^2=g_x\alpha(g_x)\in L_x(F).
\end{equation}
Changing $g_x$ multiplies $\tau_x$ by an element of $L_x(F)$,
so its action on invariant distributions is independent of this choice.
Since $\tau_x^2\in L_x(F)$, the induced action on
$L_x(F)$-invariant distributions on $N_x$ is an involution. Denote the
restricted sign character by $\varepsilon_x$.

\begin{corollary}\label{m1:cor:HC-invariance}
Under these hypotheses, if
$\mathcal D(N_x)^{\widetilde L_x,\varepsilon_x}=0$ at every closed rational
orbit, then every $L(F)$-invariant distribution on $X(F)$ is
$\tau$-invariant.
\end{corollary}
\begin{proof}
We apply Theorem~\ref{m1:thm:HC-descent} to the reductive algebraic finite
extension $\mathbf{\widetilde L}=L\rtimes_\alpha C_2$. Over $\Fb$ an
extended orbit is a union of at most two $L$-orbits of equal dimension.
If that union is closed, the boundary of either constituent lies in the
union. Boundary orbits have strictly smaller dimension, so this boundary
is empty. Conversely, a finite union of closed constituent orbits is
closed. The geometric--analytic comparison for both groups therefore
shows that the finite extension has exactly the same semisimple points.
Its orbit tangent equals the ordinary orbit tangent, since the extension
is finite. Hence its normal space is $N_x$, and its stabilizer and character
are precisely \eqref{m1:eq:corrected-extension}. The descent theorem applies.
\end{proof}

\subsection{Frobenius descent and orbital measures}\label{m1:subsec:Frobenius}
For a locally compact group $M$ our modular convention is
\begin{equation}\label{m1:eq:modulus}
 \int_M f(mh)\,dm=\Delta_M(h)\int_M f(m)\,dm,
 \qquad \mu_{L,H}=\Delta_L|_H\Delta_H^{-1}.
\end{equation}
For an algebraic group over $F$, $\Delta_M(h)$ is
$|\det\operatorname{Ad}(h)|$ on its Lie algebra.

\begin{proposition}[Frobenius descent {\cite[\S1.5]{Bernstein84}}]
\label{m1:prop:Frobenius-descent}\label{m2:tools:frobenius}
Let $p:X\to Y$ be a continuous equivariant map of $l$-spaces, where $L$
acts transitively on $Y\simeq L/L_{y_0}$ with its quotient topology. For
any character $\chi$ of $L$ there is a natural isomorphism
\begin{equation}\label{m2:tools:frobformula}
 \mathcal D(X)^{L,\chi}\simeq
 \mathcal D(p^{-1}(y_0))^{L_{y_0},\,\chi|_{L_{y_0}}\mu_{L,L_{y_0}}}.
\end{equation}
\end{proposition}
All homogeneous bases used below are rational algebraic orbits with this
topology. At $y_0$ the base density transforms by
$\Delta_{L_{y_0}}/\Delta_L=\mu_{L,L_{y_0}}^{-1}$, explaining the direction
of the fiber character. For example, projecting $F^2\setminus\{0\}$ to
$\mathbf P^1(F)$ gives the upper triangular stabilizer $P$ at $Fe_1$,
with $\Delta_P(\diag(\alpha,\beta))=|\alpha/\beta|$. Additive measure has
character $|\alpha\beta|^{-1}$, and its radial measure $|c|\,dc$ has
character $|\alpha|^{-2}=|\alpha\beta|^{-1}|\beta/\alpha|$.

\begin{lemma}[Homogeneous spaces and positive measures]
\label{m1:lem:orbital-measures}\label{m2:tools:orbit}\label{m1:cor:homogeneous-transpose}
For $X=L/H$, the space $\mathcal D(X)^{L,\chi}$ is zero unless
$\chi|_H\mu_{L,H}=1$, and is one-dimensional when this condition holds.
If $\chi$ is positive, the latter space has a positive measure generator.
Suppose an operator $\sigma$ preserves this space and the orbit, and its
square acts as the identity on the space. Then $\sigma$ fixes that
positive generator. In particular there is no distribution in the space
on which $\sigma$ acts by $-1$.
\end{lemma}
\begin{proof}
Frobenius descent reduces the assertion to a one-point fiber, whose
distribution space is one-dimensional with the stated character condition. When $\chi$ is positive, the
usual Haar-density construction gives a positive equivariant measure.
Pushforward by $\sigma$ acts on this ray by a positive scalar. Its square
is one, so that scalar is one.
\end{proof}
For $\sigma^2=h\in L$, the square condition is equivalent to $\chi(h)=1$.

\begin{proposition}[Descent with an involution]\label{m1:prop:odd-Frobenius}
Suppose $p:X\to Y$ is $\widetilde L$-equivariant and $L$ is transitive on
$Y$. Choose $g_0\tau(y_0)=y_0$, put $\tau_{y_0}=g_0\tau$, and let
$\widetilde L_{y_0}=L_{y_0}\sqcup L_{y_0}\tau_{y_0}$. Then
\begin{equation}\label{m1:eq:odd-Frobenius}
 \mathcal D(X)^{\widetilde L,\varepsilon}
 \simeq\mathcal D(p^{-1}(y_0))^{\widetilde L_{y_0},\,
 \varepsilon|_{\widetilde L_{y_0}}\widetilde\mu_{y_0}},\qquad
 \widetilde\mu_{y_0}=
 \Delta_{\widetilde L}|_{\widetilde L_{y_0}}
 \Delta_{\widetilde L_{y_0}}^{-1}.
\end{equation}
\end{proposition}
Proposition~\ref{m1:prop:Frobenius-descent}, applied to the extended
group, gives the claimed identification, with the relative modulus
evaluated on both components. An order-two
correcting operator has modulus one. In the non-unimodular applications below we check that the induced operator
preserves the positive ordinary fiber character, as required by
Lemma~\ref{m1:lem:orbital-measures}.

\subsection{Localization and opposite weights}\label{m1:subsec:localization}
\begin{proposition}[Bernstein localization {\cite[\S1.4]{Bernstein84}}]
\label{m1:prp:localization}\label{m2:tools:localization}
Let an $l$-group $L$ act on an $l$-space $X$, and let $p:X\to Y$ be a
continuous map to a Hausdorff $l$-space, constant on $L$-orbits. If
$\mathcal D(p^{-1}(y))^{L,\chi}=0$ for every $y\in Y$, then
$\mathcal D(X)^{L,\chi}=0$.
\end{proposition}
For the sign character, $p$ is required to be invariant under the extended group.

\begin{lemma}[Hyperbolas and the two-axis node]
\label{m1:lem:nonzero-xy-fiber}\label{m1:lem:Prasad-two-axis}\label{m1:subsec:two-axis}
For $c\in F$, let $V_c=\{(x,y):xy=c\}$, with
$t(x,y)=(tx,t^{-1}y)$ and $\sigma(x,y)=(y,x)$. Every
$F^\times$-invariant distribution on $V_c$ is $\sigma$-invariant.
\end{lemma}
\begin{proof}
For $c\neq0$, the equivariant parametrization
$a\mapsto(a,c/a)$ identifies $V_c$ with $F^\times$ and $\sigma$ with
$a\mapsto c/a$. Multiplicative Haar measure is invariant under translation
and inversion. For $c=0$, consider the closed fiber
\[
 Z=V_0=(F\times\{0\})\cup(\{0\}\times F).
\]
Let $T\in\mathcal D(Z)^{F^\times}$. Its restriction to the open
subset $Z\setminus\{0\}$ is
$c_1\nu_1+c_2\nu_2$, where $\nu_i$ is multiplicative Haar
measure on the corresponding punctured axis, normalized so that
the units have measure one. This follows from Frobenius descent
on each of its two open orbits, each with trivial stabilizer.
Let $\mathfrak o$ be the ring of integers and $\varpi$ a
uniformizer. The set $C=Z\cap\mathfrak o^2$ is compact open
in $Z$. For $t=\varpi$, the first-axis part of $tC$ is
$\varpi\mathfrak o$, and its second-axis part is
$\varpi^{-1}\mathfrak o$. Therefore
\[
 \mathbf1_C-\mathbf1_{tC}
 =\mathbf1_{\mathfrak o^\times\times\{0\}}
  -\mathbf1_{\{0\}\times\varpi^{-1}\mathfrak o^\times}.
\]
The difference vanishes near the origin. Invariance and the
normalization of the two Haar measures give
$0=T(\mathbf1_C-\mathbf1_{tC})=c_1-c_2$.
Thus $T-\sigma_*T$ is supported at the origin. By
\cref{m2:tools:support} it is a multiple of $\delta_0$; it is also
anti-invariant under $\sigma$, whereas $\delta_0$ is fixed.
It vanishes; see also \cite[Lemma~4.6]{Prasad90}.

\end{proof}

\begin{proposition}[Opposite-weight representations]
\label{m1:prop:opposite-weight-plane}\label{m2:elem:split}\label{m1:cor:opposite-weight-character}
For the same action and swap, every invariant distribution on $F^2$ is
fixed by $\sigma$. The assertion remains true when the action is pulled
back along a surjective character $\alpha:S\to F^\times$.
\end{proposition}
\begin{proof}
Proposition~\ref{m1:prp:localization}, applied to $(x,y)\mapsto xy$,
and Lemma~\ref{m1:lem:nonzero-xy-fiber} give the conclusion. Surjectivity makes $S$-invariance exactly
the same as $F^\times$-invariance.
\end{proof}

\begin{corollary}[Branch coefficients]\label{m1:cor:two-axis-coefficients}
Let $\nu_x,\nu_y$ be positive invariant measures on the punctured
branches of $V_0$, normalized by $\sigma_*\nu_x=\nu_y$. If an invariant distribution
on the node has restrictions $c_x\nu_x,c_y\nu_y$, then $c_x=c_y$.
\end{corollary}
\begin{proof}
The equality $\sigma D=D$ from Lemma~\ref{m1:lem:Prasad-two-axis},
restricted to the open punctured $y$-branch, gives $c_x=c_y$.
\end{proof}

\begin{lemma}[Fixed parameters and independent products]
\label{m1:lem:passive-products}\label{m2:tools:passive}\label{m1:cor:opposite-weight-passive}
Suppose every $S$-invariant distribution on $V$ is fixed by $\sigma$.
The same holds on $V\times W$ for the product actions
$s(v,w)=(sv,w)$ and $\sigma(v,w)=(\sigma v,w)$.
If $(\alpha_1,\ldots,\alpha_m):S\to(F^\times)^m$ is jointly
surjective, the independent opposite-weight action on a product of
planes, hyperbolas, or two-axis nodes fixes every invariant distribution
under each individual coordinate swap, and hence under their product.
\end{lemma}
\begin{proof}
Fixing a test function in $W$ gives an invariant distribution on $V$.
Compactly supported locally constant functions on a product are finite
sums of product test functions: partition a compact support into finitely
many compact open rectangles on which the function is constant.
This proves the first assertion. For the second, joint surjectivity
provides each factor's scaling group while fixing all other factors.
The result follows by applying the first assertion successively to
each factor, together with Lemma~\ref{m1:lem:nonzero-xy-fiber} or
Proposition~\ref{m1:prop:opposite-weight-plane} as appropriate.
\end{proof}

\subsection{Quadratic norm fibers and matrix pairs}
\begin{lemma}[Quadratic norm fibers]\label{m2:elem:norm}
Let $L/F$ be a quadratic field extension, with nontrivial
automorphism $w\mapsto\bar w$, and let
$L^1=\{q\in L^\times:N_{L/F}(q)=1\}$. For the multiplication
action of $L^1$ on $L$, every invariant distribution is fixed by
conjugation.
\end{lemma}

\begin{proof}
The $L^1$-orbits lie in fibers of the norm map $N_{L/F}:L\to F$.
A nonempty nonzero fiber
is $w_0L^1$ and is simply transitive under $L^1$. In the orbit
coordinate $w_0q$, conjugation sends
\[
 q\longmapsto(\bar w_0/w_0)q^{-1}.
\]
The multiplier lies in $L^1$, and both inversion and translation
preserve Haar measure on this abelian group. The stabilizer is
trivial, so its modulus and that of $L^1$ are both one. The zero
fiber is $\{0\}$ because $L$ is a field; its Dirac
distribution is fixed. All other fibers are empty, so
\cref{m2:tools:localization} completes the proof.
\end{proof}

\begin{lemma}[Independent norm coordinates]\label{m2:elem:products}
Let $E/F$ be a quadratic field extension, and let $(E^1)^r$ act
coordinatewise on $E^r$. Let the odd
operator be coordinatewise conjugation. Then its sign
distribution space is zero. The same conclusion holds after
pulling back the action through a surjective homomorphism
$S(F)\to(E^1)^r$, and after adjoining any passive parameter
space fixed by the extended group.
\end{lemma}

\begin{proof}
By \cref{m2:tools:passive}, the other coordinates may be treated as
fixed parameters when applying \cref{m2:elem:norm} to each coordinate. Each
coordinate conjugation fixes every invariant distribution, so their
product does also. Pullback through a surjective homomorphism gives
the same transformations and invariant distributions. The last
assertion follows again from \cref{m2:tools:passive}.
\end{proof}

\begin{lemma}[Simultaneous conjugation of matrix pairs]
\label{m1:lem:matrix-pair-transpose}
Every distribution on $M_2(F)^2$ invariant under simultaneous
$\GL_2(F)$-conjugation is fixed by componentwise transpose.
\end{lemma}
\begin{proof}
For $2\times2$ matrices put
$\bar P=(\tr P)I-P$ and
\[
 J_0=\begin{pmatrix}0&1\\-1&0\end{pmatrix},\qquad
 \bar P=J_0P^tJ_0^{-1},\qquad
 \jmath(P,Q)=(\bar P,\bar Q).
\]
For $s,t\in F$ the invariant open sets
\begin{equation}\label{m1:eq:scalar-translation-cover}
 U_{s,t}=\{(P,Q):P+sI,\ Q+tI\in\GL_2(F)\}
\end{equation}
cover $M_2(F)^2$, since each determinant excludes at most two scalar
shifts and $F$ is infinite. Translation
$T_{s,t}(P,Q)=(P+sI,Q+tI)$ identifies $U_{s,t}$ equivariantly with
$\GL_2(F)^2$ and commutes with $\jmath$, because
$\overline{P+sI}=\bar P+sI$. These opens are also $\jmath$-stable.
Prasad's simultaneous-conjugation theorem
\cite[Proposition~4.5]{Prasad90} applies on the invertible pair space and
gives $(\jmath D-D)|_{U_{s,t}}=0$ for every invariant distribution $D$.
The open cover gives $\jmath D=D$ on all matrix pairs. Since $\jmath$
is componentwise transpose followed by simultaneous conjugation by
$J_0$, this is the required transpose invariance.
\end{proof}

\subsection{Fourier support and homogeneity}
Fix a nontrivial continuous unitary additive character $\psi:F\to\mathbb C^\times$.
For a nondegenerate symmetric bilinear form $B_V$ on a finite
dimensional $F$-space $V$, use its self-dual Haar measure to define
\[
 (\mathcal F_V f)(v)=\int_V f(w)\psi(B_V(v,w))\,dw.
\]
The same formula in the $V$ variable defines partial Fourier
transform on $V\times M$. On distributions, set
$(\mathcal F_VT)(f)=T(\mathcal F_Vf)$.

\begin{theorem}[Odd-dimensional homogeneity]\label{m2:tools:homogeneity}
Let $V$ have odd dimension and nondegenerate symmetric form $B_V$,
and let $M$ be an analytic manifold over $F$. Put
$\mathcal N_V=\{v:B_V(v,v)=0\}$. A distribution on $V\times M$
whose support and partial Fourier support are both contained in
$\mathcal N_V\times M$ is zero.
\end{theorem}

This is the odd-dimensional conclusion of
\cite[Theorem~5.1.5]{AGDescent}, including its analytic parameter $M$.
In the non-Archimedean category its Schwartz distributions are the
distributions used here. In both identity descendants, $\dim_FV=3$.
